\subsection{The uniform coefficient-three obstruction on every elliptic cubic twist}

\begin{statement}{A.70}{The first finite congruence obstruction uniformly across cubic twists}{Proved}
Let $E_d$ be the compact cubic
$X^3+Y^3=dZ^3$, $d\in\Q^\times$, with origin
$O=(1:-1:0)$, and let $R\in E_d(\Q)$ be nontorsion.
Remove $D=E_d[3]$ and use the actual joint compact and
j-line level-three quotient $\Gamma_3$ of A.67, with
$J=X^3/(dZ^3)$.
If $a\in\widehat\Z$ agrees with an integer $n$ at
every coefficient prime other than three, then
$a\kappa(R)$ has a $\Gamma_3$ Selmer lift if and only
if $a=n$. This includes $n=0$ and twists ramified at two.

Combining this assertion with A.68, for every such
twist and every nonzero integer $n$, the same exact
classification holds among \emph{all} finite alterations
of $n$. Thus every nontrivial one-prime alteration
of a nonzero integer dies already at the first listed
finite congruence refinement, uniformly across these
twists. Every diagonal cubic with a rational infinity
cusp can be put in this form by rational rescaling.
The rational-origin hypothesis is explicit; no
assertion about a torsor without such an origin,
arbitrary infinite alterations, or finite-quotient
fibres is made.
\end{statement}
\begin{proof}
Take $L=\Q(E_d[3])$, and choose a normalized flex
ratio $g$ with divisor $3(d_1)-3(d_0)$ for two
cusps in the $J=1$ block, with $g(O)=1$.
Its quadratic character factors through
$\Gamma_{3,L}$ by A.67. Write $B=nR$ and
$\delta=a_3-n$, and suppose $\delta\ne0$.
The proper three-adic torsion and $R$-division
tower over $L$ is pro-three. Simultaneous
zero-primary Frobenius, quadratic disjointness,
and Chebotarev therefore force the global
$g$-class to be $\kappa_2(g(B))$ when $n\ne0$.
When $n=0$ they force the trivial class:
$g$ extends as a unit at $O$, has value one,
and kills that cusp's inertia. These are the
same forcing arguments proved in A.68.

For every three-primary torsion point $S$
with $B+S$ interior, we consequently have
\[
 \frac{g(B+S)}{g(B)}\quad\hbox{a square in the
 appropriate finite three-torsion field}.
 \tag{A70.1}
\]
For $B=O$ interpret the denominator as one.
To justify the assertion, choose a three-primary
point $T$ with $\delta T=S$. If the ratio were
nonsquare, it would remain so after the odd
extension defining $T$. Its quadratic extension
would be disjoint from the pro-three division
tower. The prescribed-primary lemma in A.68
then permits $R_3=T$ and nontrivial quadratic
Frobenius simultaneously. The actual compact
local centre is $B+S$, contradicting the forced
global class. Thus (A70.1) holds at every
required torsion level. For $n\ne0$ all
$B+S$ are interior. For $n=0$ we use only
$S\in E_d[9]\setminus E_d[3]$ below.

\par\smallskip\noindent\textbf{The dyadic cubic isomorphism and its fields.}\enspace
Let $E_0:y^2+y=x^3-1$ be the least curve of
A.69, also written as the cubic with right-hand
coefficient nine. Choose $t^3=9/d$.
The map $(X:Y:Z)\mapsto(tX:tY:Z)$ identifies
$E_d$ with that Fermat model of $E_0$,
preserves $J$ and $O$, and identifies the
$J=1$ cusp blocks. Choose in $E_d[3]$ the
preimages of the two specified flexes on
$E_0$. Their divisor ratio normalized to
value one at $O$ is $L$-rational: its
divisor is $L$-rational, and normalization
removes the multiplicative descent
ambiguity. The normalized transported
unit is exactly
\[
 g_0=\frac{b^2\zeta^2x-y-2}{b^2x-y-2},
 \qquad b^3=3,\quad\zeta^2+\zeta+1=0,
 \tag{A70.2}
\]
because its divisor agrees and its value
at $O$ is one. This identifies the actual
retained character, including its constant,
even when the global cube root defining
the isomorphism does not itself lie in $L$.

Every unit of $\Q_2$ is a cube, by Hensel's
lemma. Put $K_2=\Q_2(\zeta)$ and $F=K_2(t)$.
Either the isomorphism is defined over
$\Q_2$ and $F=K_2$, or $F/K_2$ is cyclic,
totally ramified, of degree three. In the
second case a generator $\tau$ changes the
isomorphism by a nonidentity CM automorphism.
For the image $B'$ of $B$ this gives
\[
 \tau B'=[\zeta]^kB',\qquad k\in\{1,2\}.
 \tag{A70.3}
\]
All of $E_0[3]$ is rational over $K_2$:
the odd cube root of three lies in
$\Q_2$. Thus the local three-torsion
field of $E_d$ is exactly $F$.
In the nontrivial case, faithfulness follows
because the order-three CM automorphism
acts nontrivially on $E_0[3]$.
Higher three-torsion fields have the form
$FM_r$ with $M_r/K_2$ unramified, since
$E_0$ has good reduction at two. Hence
(A70.1) holds in these local fields after
the stated isomorphism; no uncontrolled
even-degree local extension has been added.

\par\smallskip\noindent\textbf{Move the centre into the formal group.}\enspace
If $F=K_2$, then $B'\in E_0(\Q_2)$.
Its reduction lies in $E_0(\mathbf F_2)$,
which has order three. Let $B_0$ be the
unique three-torsion lift of that reduction.
It belongs to $E_0[3](\Q_2)$.
If $F/K_2$ is ramified, inertia acts
trivially on reduction, so (A70.3) shows
that $\overline{B'}$ belongs to
$\ker([\zeta]-1)$, an order-three
subgroup. Its unique three-torsion lift
$B_0\in E_0[3](K_2)$ is fixed by
$[\zeta]$ and by $\tau$.
In either case
\[
 C=B'-B_0
\]
belongs to the formal group. It is nonzero
when $n\ne0$, since $B$ is nontorsion.
For $n=0$ set $C=O$.
Dividing the relation (A70.1) at $S-B_0$
by the one at $-B_0$ gives
\[
 r_T:=\frac{g_0(C+T)}{g_0(C)}
          \quad\hbox{a square in }FM_r
 \tag{A70.4}
\]
for every three-primary torsion target
used below. If $C=O$, this is the direct
relation with denominator one. In the
ramified case $\tau C=[\zeta]^k C$.

The standard integral formal parameter
$z=-x/y$ satisfies the exact identity
\[
 z([\zeta]Q)=\zeta z(Q),
 \tag{A70.5}
\]
since $[\zeta](x,y)=(\zeta x,y)$.
For $C\ne O$ let $h=v_F(z(C))\ge1$,
using normalized integer valuation.

\par\smallskip\noindent\textbf{Odd formal height.}\enspace
If $h$ is odd, take $T=d_0$, the pole
of $g_0$. Translation by this torsion
point is an automorphism of the good
integral model. The pole has order
three and unit leading coefficient:
its denominator reduces to $x+y$,
whose intersection with the reduced
curve is $(x+1)^3$, while the numerator
is nonzero at that point. Therefore
\[
 v_F(g_0(C+d_0))=-3h,\qquad
 v_F(g_0(C))=0.
\]
The ratio has odd valuation, which
remains odd in every unramified
extension $FM_r/F$. It cannot be a
square, contradicting (A70.4).

\par\smallskip\noindent\textbf{Even height over the unramified field.}\enspace
Suppose $F=K_2$ and $h$ is even,
hence at least two. Then $C$ is
$O$ modulo four in the formal group
and $g_0(C)\equiv1\pmod4$.
For every $T\in E_0[9]\setminus E_0[3]$
the unit $g_0$ is regular and nonzero
at its reduction, so integral
translation gives
\[
 r_T\equiv g_0(T)\pmod4.
\]
Arithmetic Frobenius $\sigma_4$ fixes
$C$ and acts on $T$ by $[-2]$.
For a square unit $u^2$ in an
unramified dyadic field,
$\sigma_4(u^2)\equiv(u^2)^4\pmod4$.
Using (A70.4), we obtain
\[
 g_0([-2]T)\equiv g_0(T)^4\pmod4.
 \tag{A70.6}
\]
The same argument applies directly
when $C=O$. By the exact calculation
(A69.11), (A70.6) makes every such
$\overline T$ a zero of the nonzero
function with pole order $23$.
But $E_0[9]\setminus E_0[3]$ has
$81-9=72$ points with distinct
reductions. This is impossible.

\par\smallskip\noindent\textbf{Even height over the ramified cubic.}\enspace
Finally suppose $[F:K_2]=3$ and
$h$ is even. From (A70.3)--(A70.5)
we have
\[
 z(C)^3\in K_2,\qquad
 v_2(z(C)^3)=h\ge2.
 \tag{A70.7}
\]
Fix an unramified extension $M/K_2$
containing the required torsion.
It is disjoint from $F$.
For $T\in E_0[9]\setminus E_0[3]$,
the expression $g_0(T+Q(z))$ is an
integral power series over
$\mathcal O_M$, with constant term
$g_0(T)$. This follows from good
integral translation and the fact
that $g_0$ is a unit at
$\overline T$. Its product over
$z,\zeta z,\zeta^2z$ is invariant
under $z\mapsto\zeta z$, hence is
an integral power series in $z^3$.
The constant term is $g_0(T)^3$.
Consequently (A70.7) gives
\[
 N(T):=\operatorname{Norm}_{FM/M}
              g_0(C+T)
       =\prod_{i=0}^2g_0(T+[\zeta]^i C)
       \equiv g_0(T)^3\pmod4.
 \tag{A70.8}
\]
The denominator norm is exactly
\[
 \prod_{i=0}^2 g_0([\zeta]^i C)=1,
 \tag{A70.9}
\]
because the three line ratios
telescope. For $C=O$ the same
identity follows from $g_0(O)=1$.
Thus (A70.4) makes $N(T)$ a
square unit of the unramified
field $M$.

The rational function $N$ has
coefficients in $K_2$: the cubic
Galois action permutes its three
factors. Hence $\sigma_4 N(T)
=N([-2]T)$. Taking the square-unit
congruence and applying (A70.8)
yields
\[
 g_0([-2]T)^3\equiv g_0(T)^{12}\pmod4.
\]
Modulo two, $g_0([-2]T)=g_0(T)^4$.
The other factor in the difference
of cubes therefore reduces to
the nonzero square of $g_0(T)^4$.
It is a unit, and division gives
(A70.6). Again the $72$ reductions
would all be zeros of the function
of pole order $23$ from A.69,
a contradiction. When $C=O$ the
same norm argument applies with
$N(T)=g_0(T)^3$, so no value at
a zero or pole cusp is used.

Every case contradicts $\delta\ne0$.
If $a=n\ne0$, the point $nR$ is
interior and supplies the lift;
if $a=n=0$, the rational $O$-cusp
supplies it. This proves the pure
three-component assertion.
For a finite alteration of a
nonzero integer, A.68 already
excludes any changed prime other
than three; the assertion just
proved excludes the remaining
case. These comparisons concern
the actual retained finite
congruence character throughout.
\end{proof}
